\section{Detection by degree truncations}\label{sec:detection}

We now show that degree truncations detect every element of $\St(R)$.
The central kernel is the essential case. For a representative whose
entries are sufficiently long, phases turn cyclic-shift evaluation
into an odd power of the original element. For an element invisible
in all truncations, Lemma~\ref{lem:cyclic-vanishing} makes the same
evaluation trivial. Its already established $2$-power order then
forces the element itself to be trivial.

Keep $L,N,M$ as in Section~\ref{sec:long-entries}. For a root word
with entries of lengths at least $M$, the following lemma puts the
start and end phases directly into the coordinate labels.

\begin{lemma}[Relocation by phases]\label{lem:relocation}
Suppose $z\in J(R)$ is represented by a finite word
\begin{equation}\label{eq:high-word}
 z=\prod_{\nu=1}^m x_{i_\nu j_\nu}(w_\nu),
 \qquad |w_\nu|\ge M,
\end{equation}
where every $w_\nu$ is a nonzero word class. Label new coordinates
by $(i,s)$, with $s\in\Z/L$. In the stable group,
\begin{equation}\label{eq:relocation}
 z=\prod_{\nu=1}^m
 x_{(i_\nu,\sigma(w_\nu)),(j_\nu,\tau(w_\nu))}(w_\nu).
\end{equation}
\end{lemma}

\begin{proof}
The empty word presents no issue, so assume $m\ge1$, and let $S$
contain the finitely many original indices and at least three
coordinates. Put
$A=\F_2\oplus R_{\ge M}\subseteq R$.
We shall compare several evaluations of the word
\eqref{eq:high-word}. Choose numerical columns
$u_s\in\F_2^L$ and rows $v_s\in(\F_2^L)^*$ such that
$v_su_s=1$ for every $s\in\Z/L$. On the word basis of the tail,
set
\begin{equation}\label{eq:phase-rank-one}
 \rho(w)=u_{\sigma(w)}v_{\tau(w)}\,w,
 \qquad \rho(1)=\Id_L,
\end{equation}
and extend linearly to $A$.
This is a unital algebra homomorphism $A\to\Mat_L(R)$.
Indeed, a nonzero product $ww'$ has matching middle phases and
inherits the outside phases, so
\[
 \rho(w)\rho(w')
 =u_{\sigma(w)}
   \bigl(v_{\tau(w)}u_{\sigma(w')}\bigr)
   v_{\tau(w')}ww'
 =\rho(ww').
\]
If $ww'=0$, the product of the two entry matrices is zero as well,
regardless of the middle numerical pairing. Scalar products are
immediate.

Applying $\rho$ to the entries in \eqref{eq:high-word}, and then
using the block roots of Lemma~\ref{lem:block-laws}, produces an
element of $J(R)$. To justify this assertion, the original matrix
word is the identity already over the subalgebra $A$, since all its
entries belong to $A$ and $A\hookrightarrow R$ is injective. Its
image under $\rho$ is therefore the identity matrix. This uses only
the matrix identity; no injectivity assertion about
$\St(A)\to\St(R)$ is needed.

Write $\mathbf e_s$ for the standard numerical columns and
$\mathbf e_s^*$ for their rows. Start with
$u_s=\mathbf e_0$ and $v_s=\mathbf e_0^*$ for every $s$.
The evaluated word is $z$ supported on the coordinates $(i,0)$,
which is the same stable element by
Lemma~\ref{lem:permutations} and centrality. We change this
evaluation in three steps.

\emph{Add entries in unused columns.}
Keep all $u_s=\mathbf e_0$, and change the rows to
\[
 v_0=\mathbf e_0^*,\qquad
 v_t=\mathbf e_0^*+\mathbf e_t^*\quad(t\ne0).
\]
Every pairing $v_tu_s$ is still $1$.
Each block root acquires at most one additional factor, going
from a coordinate $(i,0)$ to an unused coordinate $(j,t)$ with
$t\ne0$. These factors belong to the strip from
$S\times\{0\}$ to $S\times(\Z/L\setminus\{0\})$.
The original factors, supported on $S\times\{0\}$, normalize this
strip by Lemma~\ref{lem:strip}. We can therefore collect the
additional factors into a single strip element. Both the old and
new evaluated words have identity matrix image, so that strip
element has identity matrix image. Strip injectivity makes it
trivial. The new evaluation is still $z$.

\emph{Change the numerical basis.}
The rows $v_t$ just defined form a basis. Let $Q$ be the numerical
matrix with these rows, so that $v_tQ^{-1}=\mathbf e_t^*$.
Also $Q\mathbf e_0=\sum_s\mathbf e_s$, because
$v_t\mathbf e_0=1$ for every $t$. Conjugating all entry matrices
by $Q$ therefore changes the columns and rows to
\[
 u_s=\sum_{r\in\Z/L}\mathbf e_r\quad\text{for every }s,
 \qquad v_t=\mathbf e_t^*.
\]
Apply the exact block-conjugation assertion of
Lemma~\ref{lem:block-laws} with $P=Q^{-1}$. This simultaneous
basis change is induced by conjugation in $\St(R)$.
The evaluated word belongs to
the central kernel, so it remains $z$.

\emph{Remove the factors in the wrong rows.}
Give coordinate $(i,s)$ the designated phase $P((i,s))=s$.
A block root with entry $w$ now factors as
\[
 \prod_{r\in\Z/L}x_{(i,r),(j,\tau(w))}(w).
\]
Its factors commute. The factor with $r=\sigma(w)$ is the root
in \eqref{eq:relocation}; all the other factors are wrong-row
generators of $U(S\times\Z/L,P)$.
Lemma~\ref{lem:wrong-row} says that the retained, consistent roots
normalize this subgroup. Collecting the other factors expresses
the evaluated word as the right side of \eqref{eq:relocation}
times an element $u\in U(S\times\Z/L,P)$.
The retained word itself has identity matrix image: apply
\eqref{eq:phase-rank-one} with
$u_s=\mathbf e_s$ and $v_t=\mathbf e_t^*$.
Hence $u$ has identity matrix image. Injectivity on the wrong-row
subgroup gives $u=1$, proving \eqref{eq:relocation}.
\end{proof}

\begin{proposition}\label{prop:odd-power}
Under the hypotheses of Lemma~\ref{lem:relocation},
\begin{equation}\label{eq:odd-power}
 W_z(C_L)=z^L.
\end{equation}
\end{proposition}

\begin{proof}
For a homogeneous word $w$ of degree $d=|w|$, the cyclic-shift
evaluation of $x_{ij}(w)$ is the block root with entry $wC_L^d$.
Since $\tau(w)-\sigma(w)=d$ in $\Z/L$, its factors are
\begin{equation}\label{eq:offset-factors}
 x_{(i,\sigma(w)+a),(j,\tau(w)+a)}(w),
 \qquad a\in\Z/L.
\end{equation}
We claim that any two such factors arising from different offsets
$a\ne a'$ commute, even when they come from different roots of the
original word.

Write their entries as $w,w'$. A match between the column coordinate
of the first factor and the row coordinate of the second requires
$\tau(w)+a=\sigma(w')+a'$. Since $a\ne a'$, this implies
$\tau(w)\ne\sigma(w')$, and hence $ww'=0$.
A match in the reverse direction similarly implies $w'w=0$.
Thus the Steinberg relations prove the claim unless the two roots
are opposite. In that case both products vanish, and the reverse
match also gives $\sigma(w)\ne\tau(w')$.
Lemma~\ref{lem:long-commutator}, whose length hypotheses hold
because $|w|,|w'|\ge M\ge N$, then gives
$[x_{uv}(w),x_{vu}(w')]=1$. This proves the claim in every case.

We may consequently collect the product by offsets, preserving
within each offset the order of the original word. The offset-zero
word is precisely the relocated word of
Lemma~\ref{lem:relocation}, and every other offset word is its
coordinate permutation. They all give $z$, since $z\in J(R)$.
There are $L$ offsets, proving \eqref{eq:odd-power}.
\end{proof}

It remains to obtain a long-entry representative from triviality in
a degree truncation. Stable coordinates make this possible: a spare
coordinate separates the two factors of a long word while it is
being conjugated.

\begin{lemma}\label{lem:high-entries}
For every integer $M\ge1$, each element of
\[
 \ker\bigl(\St(R)\longrightarrow\St(R/R_{\ge2M})\bigr)
\]
can be represented by a finite product of roots whose entries are
nonzero word classes of lengths at least $M$.
\end{lemma}

\begin{proof}
For any two-sided ideal $K$ of a unital ring, the Steinberg
presentation shows that the kernel of
$\St(R)\to\St(R/K)$ is the normal closure of the roots
$x_{ij}(a)$ with $a\in K$. Indeed, after these roots are killed,
root addition makes $x_{ij}(a)$ depend only on the class of $a$
modulo $K$, and the remaining relations are exactly the
Steinberg relations over $R/K$.
For $K=R_{\ge2M}$, root addition expands every such entry in the
homogeneous word basis. It is therefore enough to handle a
conjugate $g x_{ij}(w)g^{-1}$ with $|w|\ge2M$.

Choose a representative word $w=ab$ with $|a|,|b|\ge M$.
Both parts are nonzero because $w$ is nonzero. Let $S$ contain
$i,j$ and every coordinate used by a fixed root word for $g$,
and choose $h\notin S$. The consecutive-root relation gives
\[
 g x_{ij}(w)g^{-1}
   =[g x_{ih}(a)g^{-1},\,g x_{hj}(b)g^{-1}].
\]
The two conjugated factors lie respectively in the strips from
$S$ to $\{h\}$ and from $\{h\}$ to $S$.
More explicitly, write $A_g=\pi(g)|_S$. By Lemma~\ref{lem:strip},
\begin{align*}
 g x_{ih}(a)g^{-1}
   &=\prod_{r\in S}x_{rh}\bigl((A_g)_{ri}a\bigr),\\
 g x_{hj}(b)g^{-1}
   &=\prod_{s\in S}x_{hs}\bigl(b(A_g^{-1})_{js}\bigr).
\end{align*}
Since $R_{\ge M}$ is a two-sided ideal, every displayed entry
still belongs to $R_{\ge M}$.
Expand the two strip elements into roots, expand each entry in the
word basis, and then expand their commutator. All resulting root
entries have length at least $M$. Doing this for the finitely many
conjugates in a normal-closure expression proves the lemma.
\end{proof}

\begin{theorem}\label{thm:detection}
Degree truncations detect the stable Steinberg group:
\begin{equation}\label{eq:stable-detection}
 \bigcap_{d\ge1}
 \ker\bigl(\St(R)\longrightarrow\St(R/R_{\ge d})\bigr)=1.
\end{equation}
\end{theorem}

\begin{proof}
Let $z$ lie in the displayed intersection. Its matrix image becomes
the identity modulo every $R_{\ge d}$. The grading has
$\bigcap_dR_{\ge d}=0$, so its matrix image is the identity over
$R$ and $z\in J(R)$. Taking $d=1$ also shows that $z$ maps to $1$
in $\St(\F_2)$.

By Lemma~\ref{lem:cyclic-vanishing}, $W_z(C_b)=1$ for every
sufficiently large integer $b$. Choose an odd phase modulus $L$
from Lemma~\ref{lem:phases} beyond this threshold. This choice uses
a fixed representative of $z$ and precedes every choice of a
long-entry representative. Now choose its phase threshold $N$ and
an integer $M\ge2N+k+1$.

The element $z$ is trivial in the truncation at degree $2M$.
Lemma~\ref{lem:high-entries} supplies a representative with
nonzero word entries of lengths at least $M$. Applying
Proposition~\ref{prop:odd-power} to this representative yields
\[
 z^L=W_z(C_L)=1.
\]
On the other hand, Proposition~\ref{prop:relative-torsion} gives
$z^{2^a}=1$ for some $a\ge0$, since $z\in J(R)$ has trivial
reduction to $\St(\F_2)$. The integers $L$ and $2^a$ are coprime.
Thus $z=1$.
\end{proof}
